\part{Anforderungen, Standards und Analyse}
\chapterimage{Bilder_und_Co/analyse.jpg} % Chapter heading image
%----------------------------------------------------------------------------------------
% Anforderungsfelder & Referenzrahmen
% -------------------------
\chapter{Referenzrahmen}\label{chap:anforderungen}

Im Folgenden werden welche der wichtigsten Quellen für diese Artbeit eingeordnet. 
Diese umfassen etablierte Standards aus der Automobil und Luftfahrtbranche und 
werden um Zertifizierungsempfehlungen internationaler Behörden wie EASA und 
Forschungsbeiträge bezüglich der Zertifizierung in der Luftfahrt ergänzt.

\section{Automotive-Standards}
Die Automotive-Normen liefern gereifte Prozesse und eine skalierte Sicherheitslogik; 
sie sind teils nicht KI-spezifisch, liefern aber eine belastbare KI spezifische Vergleichsfolie.

\medskip
\begin{itemize}
  \item \textbf{ISO/IEC TR 5469 \footnote{\citetitle{iso-iec-tr-5469-2024} \parencite{iso-iec-tr-5469-2024}}}\\
  Liefert eine domänenneutrale Argumentationsstruktur für KI, jedoch auf hohem Abstraktionsniveau und ohne direkte Bindung oder Domänenkonkretisierung.
  \item \textbf{ISO 26262 (Functional Safety) \footnote{\citetitle{iso-26262} \parencite{iso-26262}}} \\
  Bietet klare Rollen- und ASIL-Logik, jedoch ohne Abdeckung von ML/Daten-Themen und nur begrenzt auf die Luftfahrt übertragbar.
  \item \textbf{ISO 21448 (SOTIF) \footnote{\citetitle{iso-21448-sotif} \parencite{iso-21448-sotif}}} — nicht KI-spezifisch, aber wahrnehmungsnah.\\
  Passt mit ODD-/Szenario-Logik gut zur Abdeckungsargumentation, jedoch ohne eigenen KI-Nachweisweg und primär für Automotive gedacht.
  \item \textbf{ISO/PAS 8800 \footnote{\citetitle{iso-pas-8800-2024} \parencite{iso-pas-8800-2024}}} \\
  Versteht sich als Konkretisierung der ISO 26262 für KI. Setzt praxisnahe Leitplanken für KI, jedoch mit von der Luftfahrt abweichender Zulassungslogik.
\end{itemize}

\section{Luftfahrt-Standards}
Die Luftfahrt-Standards sind die Baseline für Bordsoftware und Werkzeuge. Sie sind 
nicht KI-spezifisch, bilden aber den verbindlichen Nachweisrahmen.

\medskip
\begin{itemize}
  \item \textbf{DO-178C (Software Assurance) \footnote{\citetitle{do-178c} \parencite{do-178c}}} \\
  Liefert auditierbare Objectives über den gesamten SW-Lebenszyklus und DAL-Logik, jedoch sind ML-spezifische Aspekte nicht adressiert
  \item \textbf{DO-330 (Tool Qualification) \footnote{\citetitle{do-330} \parencite{do-330}}} \\
  Ermöglicht proportionale Tool-Qualifizierung mit klaren Artefakten, jedoch müssen KI-spezifische Toolrisiken projektspezifisch hergeleitet werden.
\end{itemize}

\section{Luftfahrt-Zertifizierungsempfehlungen}
Diese Dokumente adressieren explizit KI in der Luftfahrt und ergänzen die Baseline um Prinzipien und Evidenzformen.

\medskip
\begin{itemize}
  \item \textbf{EASA AI Concept Paper (Issue~2) \footnote{\citetitle{easa-ai-concept-issue2-2024} \parencite{easa-ai-concept-issue2-2024}}}\\
  Gibt Leitplanken für KI und einen Level-Pfad, jedoch bleibt es high-level ohne vollständige Detailanforderungen.
  \item \textbf{EASA AI Roadmap 2.0 \footnote{\citetitle{easa-ai-roadmap-2-2023} \parencite{easa-ai-roadmap-2-2023}}} \\
  Verortet KI strategisch menschzentriert, jedoch ist sie nicht normativ und die Ableitung auf DAL-Ebene verbleibt beim Projekt.
  \item \textbf{EASA MLEAP Final Report \footnote{\citetitle{easa-mleap-final-2024} \parencite{easa-mleap-final-2024}}} \\
  Liefert konkrete Methoden und Evidenzen für KI aus luftfahrtnahen Use-Cases, jedoch als Empfehlung ohne bindenden Status und mit Fokus auf Level~1/2.
\end{itemize}

\section{Forschungsbeiträge}
Die Forschungsarbeiten schließen Lücken mit KI-fokussierten Konzepten und Checklisten für die Evidenzführung.

\medskip
\begin{itemize}
  \item \textbf{ML meets aerospace: challenges of certifying airborne AI \footnote{\citetitle{ml-meets-aerospace-2024} \parencite{ml-meets-aerospace-2024}}} — KI-bezogen.\\
  Leitet systematisch Nachweise beispielsweise für ODD-Disziplin für KIs ab, jedoch mit Überblickscharakter und wenigen domänenspezifischen KPIs.
  \item \textbf{Towards certifiable AI in aviation \footnote{\citetitle{towards-certifiable-ai-aviation-2024} \parencite{towards-certifiable-ai-aviation-2024}}} — KI-bezogen.\\
  Bietet Strukturierungshilfen für KI-Evidenzpakete, jedoch überwiegend konzeptionell mit nur punktuellen empirischen Belegen.
\end{itemize}



%---------------------------------------------
%Luftfahrt
%---------------------------------------------

\chapter{Luftfahrtrahmenbedingungen}\label{chap:anf-avi}

Folgend werden systematisch EASA Dokumente auf Forderungen an die Nachweise für KI untersucht und 
zusätzlich etablierte Normen betrachtet. 





\section{Luftfahrt-Standards \& ihre Rolle im KI-Kontext}
DO-178C bleibt die Baseline für Software-Assurance in der Avionik. Ziele und Evidenzen skalieren mit 
DAL A–E. EASA verweist explizit auf die DO-178C als Nachweisgrundlage. 
KI-spezifische Guidance kommt zusätzlich obenauf, ändert aber die Software-Pflichten nicht. 
Zudem berühren Daten-/Modell-Pipelines die Tool-Qualifikation (DO-330).

DO-178C und DO-330 beziehen sich nicht spezifisch auf KI, sondern allgemein auf Software und Tools. Deterministischer Code 
und MC/DC-Coverage sind stark empfohlen und gelten als Best Practice. Bei KI-Komponenten sind diese Punkte teils nicht erfüllbar 
und werden in den Normen nicht ausdrücklich gefordert. Daher ist KI unter diesen Standards grundsätzlich zulassungsfähig, jedoch 
steigen die Anforderungen an die Nachweisführung.






\section{EASA AI-Roadmap }\footnote{\citetitle{easa-ai-roadmap-2-2023} \parencite{easa-ai-roadmap-2-2023}}
Die EASA positioniert KI als menschzentriert und staffelt Anwendungen nach Autorität: \textbf{Level 1} (Assistenz), 
\textbf{Level 2} (Human-AI-Teaming) und später \textbf{Level 3} (fortgeschrittene Automatisierung). 
Roadmap 2.0 definiert Vertrauenswürdigkeits-Bausteine (u.\,a. Learning Assurance, Erklärbarkeit, Human Oversight) 
und skizziert einen gestuften Zulassungspfad. Betont wird, dass KI das bestehende Safety-Regelwerk ergänzt aber nicht ersetzt.
\paragraph{Was das praktisch bedeutet}
\begin{itemize}
  \item \textbf{Stufenweise Zulassung:} Erst Level-1, später Level-2 – jeweils mit klar definiertem Einsatzbereich und menschlicher Aufsicht. Kein Nachweisrabatt: Die Strenge der Nachweise richtet sich weiterhin nach Hazard/DAL.
  \item \textbf{Zusatz statt Ersatz:} KI-spezifische Evidenzen kommen oben drauf zu etablierten Software-/System-Nachweisen.
  \item \textbf{Lebenszyklus im Blick:} Betriebliches Monitoring, kontrollierte Updates und nachvollziehbare Änderungen sind Teil des Sicherheitsarguments (``living Safety Case'').
\end{itemize}



\section{EASA AI-Guidance 2024 (Concept Paper Issue 2) – Level 1 \& 2} \footnote{\citetitle{easa-ai-concept-issue2-2024} \parencite{easa-ai-concept-issue2-2024}}
Das Concept Paper beschreibt die schrittweise Einführung von KI in sicherheitsrelevanten Anwendungen: 
zuerst \textbf{Level 1} (Assistenzfunktionen), anschließend ausgewählte \textbf{Level 2}-Anwendungen (Human–AI Teaming). 
Die Guidance ergänzt das bestehende Safety-Regelwerk (z.\,B.\ DO-178C) um KI-spezifische Bausteine wie Learning Assurance, 
Erklärbarkeit, präzise Einsatzgrenzen (ODD) und betriebliches Monitoring – ohne die klassischen Nachweise zu ersetzen.

\paragraph{Kernpunkte:}
\begin{itemize}
  \item \textbf{Level-Feinbild \& „statisches“ AI-Level:} Unterscheidung L1A/L1B sowie L2A (cooperation, directed) und L2B (collaboration, supervised); das höchste erreichte Fähigkeitsniveau legt ein statisches Level fest und steuert die Proportionalität der Nachweise.
  \item \textbf{ConOps \& Aufgabenallokation (HAT):} Endnutzerzentrierte ConOps mit klaren ``task Allocation Patterns'' (Rollen, Autorität, Übersteuerbarkeit) als verbindliche Grundlage für Design, V\&V und Betrieb.
  \item \textbf{OD vs.\ ODD (Einsatzgrenzen):} Zweistufige Scope-Führung: ``Operational Domain'' auf Systemebene und verfeinertes ODD für AI/ML-Constituents; außerhalb davon gilt abstain/degradation.
  \item \textbf{Methodenumfang \& Trainingsmodus:} Fokus auf überwachtes (und initial unüberwachtes) Lernen; kein Reinforcement Learning in L1/L2; Modelle werden ``offline'' trainiert und zur Freigabe eingefroren (kein Online-Lernen im Betrieb).
  \item \textbf{Kritikalitätsgrenzen \& Assurance:} In der Anfangsphase kein Einsatz von AI/ML-Constituents in höchsten Kritikalitäten wie DAL~A/B; keine Reduktion klassischer Assurance-Anforderungen innerhalb der AI/ML-Teile.
  \item \textbf{Funktionsanalyse bis zum AI/ML-Item:} Systematische Zerlegung bis zur Trennlinie AI/ML-Item vs.\ traditionelle Komponenten, Anforderungen und V\&V zielgenau führen zu können.
\end{itemize}




\section{MLEAP Final Report (2024)} \footnote{\citetitle{easa-mleap-final-2024} \parencite{easa-mleap-final-2024}}
Der \emph{Final Report} des EASA-Forschungsprojekts MLEAP (Machine Learning Application Approval) bündelt Methoden- und 
Werkzeugempfehlungen zur Learning Assurance und zeigt, wie der EASA-W-Shape (lernprozessspezifische Ergänzung zum V-Modell) 
praktisch umgesetzt werden kann. Schwerpunkte sind Datenrepräsentativität, Generalisierung und Robustheit/Stabilität 
von ML-Modellen; ergänzt wird dies durch eine generische Entwicklungs-Pipeline im Sinne des W-Shape. Die Ergebnisse fließen 
in das \emph{EASA AI Concept Paper Issue 2} ein.

\paragraph{Kerninhalte:}
\begin{itemize}
  \item \textbf{W-Shape als Lernprozess-Rahmen:} Ergänzt das V-Modell um ML-spezifische Ziele/Verifikationen (Daten, Lernen, Modell, Integration) entlang des gesamten Zyklus.
  \item \textbf{ODD-gebundene Datenqualität:} Repräsentativität wird explizit aus dem AI/ML-ODD abgeleitet; Balance ``Completeness'' vs.\ ```Diversity'' wird methodisch geführt (Selection Grid).
  \item \textbf{``Right-Sizing'' \& Generalisierung:} Modellkomplexität wird zweckadäquat gewählt; Generalisation Guarantees werden prozessual abgesichert (Splits, Bounds, Unsicherheiten).
  \item \textbf{Stabilität vs.\ Robustheit:} Klare Trennung der Ziele und Prüfmethoden (statistisch, empirisch, formal) mit Perspektive auf Runtime-Überwachung.
  \item \textbf{Systematik seltener Fälle:} Gemeinsame Taxonomie (Edge/Corner/Outlier) und mehrlagiges Schema (Pixel→Domain→Object→Scene→Scenario) für Tests und Datenerzeugung.
  \item \textbf{Use-Case-getriebene Bewertbarkeit:} Konkrete Luftfahrt-Use-Cases dienen als Trägerartefakte für Methodenwahl, Bewertung und operativen Transfer.
\end{itemize}


\section{Forderungen aus der Luftfahrt}
Im Folgenden sind die Forderungen zur Zulassung einer KI aus der Roadmap, dem Concept Paper und dem MLEAP Report 
zusammengetragen und mit IDs zum späteren Vergleich markiert.

\paragraph {F-Avi-01 Human-centric, Level \& ConOps/HAT (Human-AI-Teaming)}
  {KI wird human-zentriert ausgelegt; statisches AI-Level (L1A/B, L2A/B) richtet die Prüftiefe aus. Eine endnutzerzentrierte ConOps mit klaren ask Allocation Patterns (Rollen, Autorität, Übersteuerbarkeit) ist verbindlich. \textit{Roadmap, Concept}}

\paragraph {F-Avi-02 OD/ODD \& Abstain/Degradation}
  {Einsatzgrenzen sind zweistufig zu führen: OD (System) und verfeinertes ODD (AI/ML-Constituent). Außerhalb des ODD verhält sich die KI konservativ (Abstain/Degradation). \textit{Roadmap, Concept}}

\paragraph {F-Avi-03 Knowledge \& Data Governance (ODD→Daten)}
  {Daten- und Wissensbasen werden vor dem Training qualifiziert (Provenienz, Repräsentativität) und explizit aus dem ODD abgeleitet; Methodenwahl wird dokumentiert (z.\,B.\ Selection Grid). \textit{Roadmap, Concept, MLEAP}}

\paragraph {F-Avi-04 AI Assurance (Learning \& Explainability)}
  {Der Übergang von Programmieren zu Lernen wird abgesichert; Erklärbarkeit ist technisch (Dev/Review) und operativ (nutzerverständlich, rechtzeitig) vorzusehen. \textit{Roadmap, Concept}}

\paragraph {F-Avi-05 Betriebs-Recording}
  {Bord-/Betriebsaufzeichnungen unterstützen kontinuierliches Sicherheitsmonitoring und Vorfallanalysen; sie sind Bestandteil des Safety Case. \textit{Roadmap}}

\paragraph {F-Avi-06 OOD \& Unvorhersagbarkeit}
  {Verteilungswechsel werden erkannt; Unsicherheit wird gemanagt und löst konservatives Verhalten aus. Runtime-Indikatoren/Monitore sind vorzusehen. \textit{Roadmap, MLEAP}}

\paragraph {F-Avi-07 Generalisierung, Robustheit \& Stabilität}
  {Generalisierung ist nachzuweisen (zeit-/raumgetrennte Splits); Stabilität (kleine Perturbationen) und Robustheit (adverse Bedingungen) werden mit getrennten Test-Suiten geprüft; ein Komplexitätsbudget verhindert Under-/Overfit. \textit{Roadmap, MLEAP}}

\paragraph {F-Avi-08 Lebenszyklus \& W-Shape}
  {Der Safety Case umfasst Aktivitäten vor dem Training und nach der Zulassung (Monitoring, kontrollierte Änderungen). Der EASA-W-Shape dient als Governance-Rahmen mit klaren Gates/Artefakten. \textit{Roadmap, MLEAP}}

\paragraph {F-Avi-09 Kein Online-Lernen im Betrieb}
  {Modelle sind offline trainiert und zur Freigabe eingefroren; Online-/Continual-Learning in der Operation ist ausgeschlossen. \textit{Roadmap, Concept}}

\paragraph {F-Avi-10 Einbettung in Safety-Systeme/Plattform}
  {Die Semantik des Modells bleibt bei Portierung/Optimierung erhalten; Plattform-/Hardware-Aspekte (inkl. Beschleuniger) sind qualifikationsfähig auszulegen. \textit{Roadmap}}

\paragraph {F-Avi-11 Rare-Events-Management}
  {Seltene/novel Situationen werden systematisch erfasst (Taxonomie, Katalog) und für Tests sowie Datenpflege genutzt. \textit{MLEAP}}

\paragraph {F-Avi-12 Use-Case-Verankerung}
  {Methoden- und Metrikentscheidungen sind an Use-Case-Anforderungen gebunden; operative Übergabeartefakte (KPIs, Monitoring-Hooks, Logging) sind vor Einsatz definiert. \textit{MLEAP}}









%--------------------------------------------------------------------------
%------------------------------------------------------------------------------
\chapter{Luftfahrt \(\leftrightarrow\) Automotive}\label{chap:automotiv}

Um zu verstehen in wiefern man die Übertragbarkeit von Automotive Standards gegenüber der Luftfahrt betrachten kann, 
hilft es zu verstehen in wiefern diese sich bezüglich KI, Sicherheitsnachweisen oder allgemein als Branche unterscheiden oder sogar ähneln. 
Was als erstes ins Auge sticht, ist eine gefahrskalierende Nachweisstrenge. In der Luftfahrt als DAL-Level und 
im Automotive als ASIL-Level bezeichnet. DAL berücksichtigt ausschließlich den Schweregrad, ASIL hingegen 
berücksichtigt zusätzlich die Eintrittswahrscheinlichkeit und die Beherrschbarkeit.
\begin{table}[h]
\centering
\caption{Gegenüberstellung: Luftfahrt (DAL) und Automotive (ASIL).}
\begin{tabular}{p{0.3\textwidth} p{0.3\textwidth}}
\toprule
\textbf{Luftfahrt (DAL)} & \textbf{Automotive (ASIL)} \\
\midrule
DAL A (catastrophic) & ASIL D (sehr hoch)\\
DAL B (hazardous) & ASIL C (hoch) \\
DAL C (major) & ASIL B (mittel)\\
DAL D (minor) & ASIL A (niedrig)  \\
\bottomrule
\end{tabular}
\end{table}

\section{KI abgeleitet nach Referenzen}
Was nach erstmaligen überfliegen der bereits genannten Referenzen deutlich wird, ist dass sowohl 
in der Luftfahrt als auch im Automotive zum Sicherheitsnachweis für KI kein Nachweisrabatt 
gewährt wird beziehungsweise die bereits etablierten Standards bezüglich KI ergänzt werden müssen. 
Zudem ist in der Luftfahrt eine Menschzentrietheit und Einstufung nach Autorität erkennbar, während 
im Automotive eine Unterscheidung nach Automatisierungsgrad gemacht wird. Der Automatisierungsgrad 
unterscheidet zwar ob Mensch oder Maschine entscheidet, setzt diesen aber nicht in den Mittelpunkt.

\section{Regulatorik}
Die ersten Unterschiede werden hingegen bereits in der Regulatorik sichtbar, während die Luftfahrt 
mit einer starken Behördenbindung eine Zulassung vor dem Betrieb vorraussetzt, kann im Automotive 
mit Typengenehmigungen und in der USA sogar mit Selbstzertifizierungen gearbeitet werden.

\section{Betrieb, Umgebung und Skalierung}
Weitere deutliche Unterschiede sind auch im Betrieb, der Umgebung und den Skalen sowohl physikalisch als auch 
in der Produktion erkennbar. Während Flugzeuge nicht einfach nur fliegen und Autos am Boden fahren sind Flugzeuge 
auch an Zeitpläne gebunden und werden oft mit einer mehrköpfigen Crew betrieben, Autos hingegen werden oft 
von Einzelpersonen überall auf der Welt spontan für Privatfahrten benutzt. Das sorgt dafür, dass Flugzeuge eher 
vorhersehbare Einsatzgrenzen haben und Autos in sehr unterschiedlichen Umgebungen benutzt werden. Was auch deutlich 
erkennbar ist, dass in der Luftfahrt man jeden Tag von einer kleien Drone bis hin zum Airbus A380 begegnenn kann, wohingegen 
im Automotive eher vom Fußgänger bis hin zum LKW ausgegangen wird. Weiter kann man im Automotive 
auf Sicht ein annäherndes Fahrzeug auch auf der Autobahn mit ein bis zweihundert km/h vorbei fahren sehen, 
in der Luftfahrt kann man ein weit entferntes Überschall Flugzeug eher nur erahnen. Was nicht sofort erkennbar ist, 
dass die meistgebaute Fluzeugserie gut 12.000 Flugzeuge umfasst und Automobilzulieferer mit Serien 
von Einhundertmillionen rechnen. 

\section{Software- und Hardware}
Für diese Arbeit wichtig ist das handhaben von Updates im Betrieb, während Automotive Beobachtungen im Feld 
und regelmäßige Updates sowie zur Not Rückrufaktionen in Betracht zieht, ist die Praxis in der Luftfahrt, 
diese in den Anforderungen und erneute Freigaben einzuplanen. 

Ähnlich ist hingegen die Kostenplanung zum 
Einsatz von Radarsystemen im allgemeinen, zur Veranschulichung eine Beispielformel:
\[\min_{\tau}\; J(\tau)
= \underbrace{c_{\mathrm{FA}}\;P_{\mathrm{FA}}(\tau)}_{\text{Kosten falscher Alarm}}
+ \underbrace{c_{\mathrm{FN}}\;P_{\mathrm{FN}}(\tau)}_{\text{Kosten False Negative}}
+ \underbrace{c_{\mathrm{sys}}(\tau)}_{\text{Systemkosten}} 
\]
Es geht darum die Kosten des Systems im Einsatz zu reduzieren. Minimiert werden soll die Summe der Kosten eines falschen Alarms, 
die Kosten etwas nicht erkannt zu haben und die Kosten das System zu verbauen. 
Zivil priorisiert das Vermeiden von Personenschäden. Im militärischen Kontext muss mit Personenschäden gerechnet werden. 
Weiter kommt hier die Frage auf, ob man die Fehlerwahrscheinlichkeit von Menschen im Vergleich  
zu einer Maschine berücksichtigen sollte.

Auch die Hardware Zulassung sollte betrachtet werden. Im Automotive werden bespielsweise GPUs aufgrund ihrer Unsicherheit 
redundant verbaut um Sicherheit zu gewährleisten. In der Luftfahrt, wo ein System ohnehin redundant verbaut wird, 
würde man diese nach dieser Vorgehensweise ``stapeln''. Eine Beispielrechnung, warum ein solcher demokratischer Entscheider 
problematisch sein kann ist: Wenn man für eine Entscheidung 11 
unabhängige Klassifizierer verbaut, welche jeweils 9 unabhängige GPUs haben, würde man eine 30 aus 99 Entscheidung treffen
(6 aus 11 notwendige Klassifizierer für eine Entscheidung, mit jeweils 5 aus 9 GPUs) 


\section{Anforderungen}
Komplexer hingegen sind zum Beispiel konkrete Zahlen in den Anforderungen. Zunächst dient eine 
Beispielrechnung der Größenordnungen zur Veranschaulichung:

\begin{table}[H]
\centering
\caption{Beispielhafte Dimensionierung für Sensoren (vereinfachtes Rechenbeispiel)}
  \begin{tabular}{p{0.40\textwidth} p{0.25\textwidth} p{0.25\textwidth}}
\toprule
\textbf{Kriterium} & \textbf{Automobilbranche} & \textbf{Luftfahrtbranche} \\
\midrule
  Betriebsstunden pro Fahrzeug (Lebensdauer) & 5.000  & 150.000 \\
  Anzahl verbauter Sensoren (Flotte) & 100.000.000 & 10.000 \\
  Gesamtbetriebsstunden (Flotte) & 5*$10^{11}$h & 15*$10^{8}$h  \\
  Zielraten (Beispiel) & $<10^{-9}$/h &  $<10^{-9}$/h \\
  Anzahl Ausfälle (Betriebsstunden * Zielrate) & < 500 &  < 1{,}5 \\
\bottomrule
\end{tabular}
\end{table}


In der Luftfahrt sind Fehlereintrittsraten nach Anwendung und Gefahreneinstufung oft eindeutig festgehalten. 
Im Automotive werden beispielsweise Fehlereintrittsraten bezügliche der Einstufung, vorallem bei niedriger Einstufung, unteranderem 
von Robustheitsanforderungen überschattet. Auch wenn theoretisch eine höhere Eintrittsrate das einzelne Fahrzeug 
nicht unangemessen unsicherer machen würde, kann dies zu Rückrufaktionen führen und somit hohe Kosten verursachen. 
Daher und zur Vermeidung zusätzlicher Redundanz werden im Automotive die Fehlereintrittsraten oft relativ streng gesetzt. 
Die Luftfahrt hat allgemein hohe Anforderungen und erhöht die Sicherheit zusätzlich durch Redundanz. 








\chapter{Deep-Dive: ISO/PAS 8800}\label{chap:iso-8800}
Die ISO 8800 \parencite{iso-pas-8800-2024} ist sehr Dokumentationslastig und beschreibt \textbf{nicht} wann eine KI sicher ist oder wie sie 
sicher wird. Diese Norm beschreibt was zu tun ist und wie man begründet, dass eine KI sicher genug ist. 
Detailprozesse und Risikobeurteilungen sind nicht enthalten. Stattdessen wie existierende Normen 
angepasst oder erweitert werden müssen.
\newcommand{\dotG}{{\color{green}\large$\bullet$}}
\newcommand{\dotY}{{\color{blue}\large$\bullet$}}
\newcommand{\dotR}{{\color{red}\large$\bullet$}}



\section{Inhalte (nach Annex A) und Bemerkungen}
\paragraph{Legende zur Einstufung der Relevanz bezüglich dieser Arbeit}
\begin{itemize}[leftmargin=2.5cm,labelsep=0.8em]
  \item[\dotG] relevant
  \item[\dotY] teilweise relevant
  \item[\dotR] nicht besonders relevant
\end{itemize}
\medskip
\begin{itemize}
  \item[\dotY] (0) \textbf{Einleitung}
  \item[\dotR] (1) \textbf{Normreferenzen}
  \item[\dotY] (2) \textbf{Scope}
  \item[\dotR] (3) \textbf{Begriffe und Definitionen}
  \item[\dotR] (4) \textbf{Abkürzungen}
  \item[\dotR] (5) \textbf{Anforderungen für Konformität:}
    Anforderungen an die ISO 26262-2 müssen erfüllbar sein oder zumindest begründete akzeptable Alternativen.
  \item[\dotY] (6) \textbf{KI im Kontext der Systemsicherheitsentwicklung für Straßenfahrzeuge und Grundkonzepte:}
    Wird berücksichtigt und im Blick auf Radarsysteme betrachtet.
  \item[\dotR] (7) \textbf{KI-Sicherheitsmanagement:} Bezieht sich auf Aktivitäten rund um den Sicherheitslebenszyklus, 
    die notwendig sind, um Sicherheit zu gewährleisten und tragen nicht zum Verständnis der technischen Sicherheit an sich bei.
    Man sollte anerkennen, dass dies wichtig ist, verlässt aber den Scope dieser Arbeit. Die Ziele sind: \\
    (a) Lebenszyklus festlegen, um unvertretbare Risiken zu vermeiden auf Fahrzeugebene;\\
    (b) geeignete organisatorische und projektspezifische Sicherheitsmanagementprozesse sicherstellen;\\
    (c) KI-Sicherheitsaktivitäten planen, initiieren und durchführen.
  \item[\dotY] (8) \textbf{Assurance-Argumente für KI-Systeme:} Die Ziele sind:\\ 
    (a) Assurance-Argument entwickeln, das die Erfüllung der zugewiesenen Sicherheitsanforderungen belegt;\\
    (b) prüfen, ob das Argument das tatsächliche Restrisiko eines Verstoßes gegen die Sicherheitsanforderungen korrekt abbildet.
  \item[\dotG] (9) \textbf{Ableitung von KI-Sicherheitsanforderungen:} Die Ziele sind:\\
    (a) Vollständige, konsistente Sicherheitsanforderungen an das KI-System spezifizieren;\\
    (b) Anforderungen auf Basis von Entwicklung/Verifikation/Validierung verfeinern;\\
    (c) Einschränkungen des KI-Systems über seinen Eingaberaum definieren und an den übergeordneten Entwicklungsprozess eskalieren;\\
    (d) erforderliche Off-board- und On-board-Maßnahmen gemäß Abschnitt 14 spezifizieren.
  \item[\dotY] (10) \textbf{Auswahl von KI-Technologien sowie Architektur- und Entwicklungsmaßnahmen:} Die Ziele sind:\\
    (a) Geeignete KI-Technologien auswählen und begründen;\\
    (b) Architektur- und Entwicklungsmaßnahmen zur Erfüllung der Sicherheitsanforderungen vor dem Einsatz identifizieren;\\
    (c) Architekturmaßnahmen zur Minderung nach dem Einsatz erkennbarer funktionaler Unzulänglichkeiten identifizieren;\\
    (d) Maßnahmen definieren, die die Erfüllung der Sicherheitsanforderungen im Zielanwendungsbereich sicherstellen.
  \item[\dotR] (11) \textbf{Datenbezogene Aspekte:} Die Ziele sind:\\
    (a) Datensatz-Lebenszyklus (Erhebung, Erstellung, Analyse, Verifikation/Validierung, Verwaltung, Wartung) definieren;\\
    (b) Datensatz-Unzulänglichkeiten mit Sicherheitsauswirkung identifizieren;\\
    (c) datenbezogene Sicherheitseigenschaften mit Einfluss auf die Systemsicherheit identifizieren und zur Analyse heranziehen;\\
    (d) Gegenmaßnahmen zur Vermeidung/Minderung von Datensatz-Unzulänglichkeiten über den Datensatz-Lebenszyklus hinweg festlegen;\\
    (e) datenbezogene Arbeitsprodukte definieren, die den Sicherheitsnachweis unterstützen.
  \item[\dotG] (12) \textbf{Verifikation und Validierung des KI-Systems:} Die Ziele sind:\\
    (a) Verifizieren, dass das KI-System seine Sicherheitsanforderungen erfüllt;\\
    (b) validieren, dass die dem KI-System zugewiesenen Sicherheitsanforderungen bei der Integration ins übergeordnete System erreicht werden.
  \item[\dotY] (13) \textbf{Sicherheitsanalyse von KI-Systemen:} Die Ziele sind:\\
    (a) Sicherheitsrelevante Faults und KI-Fehler identifizieren, die KI-Sicherheitsanforderungen verletzen können;\\
    (b) potenzielle Ursachen identifizieren;\\
    (c) Sicherheitsmaßnahmen zur Vermeidung oder Beherrschung sicherheitsrelevanter KI-Fehler definieren unterstützen;\\
    (d) Verifikation der KI-Sicherheitsanforderungen unterstützen – durch Anpassung bzw. Neuableitung von Anforderungen zu Daten-, Entwurfs- und Test­spezifikationen.
  \item[\dotR] (14) \textbf{Maßnahmen im Betrieb:} Maßnahmen im Betrieb außerhalb des Scopes dieser Arbeit.\\ 
    (a) Prozessanforderungen für die kontinuierliche Sicherstellung der KI-Sicherheit nach dem Einsatz;\\
    (b) Maßnahmen aus Abschnitt 8 und/oder zusätzliche Maßnahmen zur Risikofindung und zum Erhalt der KI-Sicherheit im Betrieb nutzen;\\
    (c) geeignete Reaktionen auf unzulässige Risiken vorsehen und die erneute Freigabe des geänderten KI-Systems vor Auslieferung sicherstellen.
  \item[\dotY] (15) \textbf{Vertrauen in die Nutzung von KI-Entwicklungsframeworks und -Softwaretools:} Die Ziele sind:
   (a) Anforderungen und Leitlinien bereitstellen, um Fehlerquellen und unangemessene Verzerrungen in Offline-Prozessen, 
   -Tools und -Prinzipien zu identifizieren, zu mindern und zu dokumentieren, die zur Entwicklung, Verifikation und 
   Bereitstellung sicherheitsrelevanter KI-Modelle genutzt werden.
  \item[\dotY] (A-H) \textbf{Anhänge (A–H):}\\
    (A){\dotR} Überblick und Ablauf dieses Dokuments\\
    (B){\dotR} Beispielstruktur eines Assurance-Arguments für ein KI-System\\
    (C){\dotG} ISO 26262-Lückenanalyse für ML\\
    (D){\dotG} Detaillierte Betrachtungen zu sicherheitsrelevanten Eigenschaften von KI-Systemen\\
    (E){\dotR} STAMP/STPA-Beispiel\\
    (F){\dotR} Identifikation von Softwareeinheiten in NN-basierten Systemen\\
    (G){\dotG} Architektur- und Entwicklungsmaßnahmen für KI-Systeme\\
    (H){\dotR} Typische Leistungskennzahlen für maschinelles Lernen\\
\end{itemize}




\section{(0) Einleitung}
Die ISO 8800 passt die ISO 26262 und ISO 21448 gezielt für KI an und ergänzt sie um ein 
kausales Modell für KI-bedingte Unzulänglichkeiten, daraus abgeleitete Sicherheitsanforderungen sowie 
Risikominderungs­maßnahmen. Weil Anwendungen und Stand der Technik stark variieren, werden weniger starre 
Prozessvorgaben gemacht und stärker auf projekt­spezifische Assurance-Argumente mit iterativer Risikoreduktion 
gesetzt. Das Dokument ist u. a. mit ISO/IEC 22989, TR 5469 abgestimmt und zeigt, wann seine Leitlinien gelten: immer 
dann, wenn Sicherheitsanforderungen entweder einer KI-Funktion oder einer KI-Sicherheitsmaßnahme zugeordnet sind.

\subsection*{Schlüsselpunkte}
\begin{itemize}
 \item Zuschnitt/Erweiterung von ISO 26262 und ISO 21448 speziell für KI.
 \item Kausales Modell für Quellen funktionaler Unzulänglichkeiten in KI; daraus abgeleitete KI-Sicherheitsanforderungen und Maßnahmen.
 \item Fokus auf projektspezifisches Assurance-Argument statt detaillierter Prozessvorgaben; iterative Risikoreduktion.
 \item Hazard/Risk-Analysen auf Fahrzeugebene sind nicht enthalten (bleiben ISO 26262/ISO 21448/ISO TS 5083 vorbehalten).
 \item Enge Anbindung an ISO/IEC 22989, TR 5469; Klassen/Usage-Levels werden nicht explizit verwendet, dienen aber der Anforderungsermittlung.
 \item Gültig für KI als Funktion oder als Sicherheitsmechanismus – auch außerhalb des unmittelbaren SOTIF-Scopes.
\end{itemize}

\section{(2) Scope}
Das Dokument gilt für sicherheitsrelevante Systeme in Serienfahrzeugen (ohne Mopeds), 
die E/E-Systeme mit KI einsetzen. Es adressiert fahrzeugseitige Sicherheitsrisiken, die durch Ausgabemängel, 
systematische Fehler und zufällige Hardwarefehler von KI-Elementen entstehen – 
inklusive Wechselwirkungen mit KI-Elementen außerhalb des Fahrzeugs, sofern diese Sicherheit direkt (z. B. 
externe Objekterkennung) oder indirekt (z. B. Felddaten-Monitoring) beeinflussen. Es beschreibt sicherheitsrelevante Eigenschaften 
von KI-Systemen zur Untermauerung eines belastbaren Safety-Assurance-Claims, gibt aber keine spezifischen Leitlinien 
für KI-basierte Softwaretools. Schwerpunktmäßig behandelt es ML-basierte Ansätze auf Prinzipienebene, ohne in methodische 
Details zu gehen (z. B. konkrete DNN-Architekturen).

\section{Kapitel Deep-Dive}
\subsection{(6) KI im Kontext der Systemsicherheitsentwicklung}
ISO/PAS~8800 ergänzt die ISO~26262, sobald ML\slash KI als Bestandteil der sicherheitsrelevanten 
Verarbeitung genutzt wird. Eine typische Architektur gliedert sich in Vorverarbeitung, 
ML-Modell (Detektion/Klassifikation/Tracking) und Nachverarbeitung/Monitoring. 
Klassische Deploy\-ment-Schritte  werden durch ISO~26262 abgedeckt; daten- und 
trainingsbezogene Schritte fallen in den Fokus der ISO/PAS~8800.

\paragraph{Interaktion mit der Systemebene}
Anforderungen und KPIs werden auf Systemebene zugewiesen und iterativ geschärft. Ist die Erfüllung 
nicht belastbar nachweisbar (fehlende Daten, Evidenz), sind Systemmaßnahmen zu wählen: 
ODD-Restriktion, Diversity (z.\,B.\ Radar+\-Kamera), Redundanz (parallelisierte Pfade). Integrationstests 
und Felddaten können weitere Iterationen anstoßen.

\paragraph{Fehlerfortpflanzung und Ursachen}
Es werden Ketten beschrieben, die Fehler anhand von Ursache oder Auswirkung beschreiben
und wie sie sich in das System eingliedern oder fortpflanzen. Beschrieben wird es nach:
KI-Triggerbedingungen, Faults bzw.\ funktionalen Unzulänglichkeiten, 
KI-Fehlern und unerwünschtem Verhalten. Radarbezogene Ursachen können sein (I) Spezifikationslücken im 
ODD, (II) Performance-Grenzen, (III) systematische Trainings-/SW-Fehler 
sowie (IV) zufällige Hardwarefehler.


\subsection{(8)Assurance-Argumente für KI-Systeme}
Ein Assurance-Argument ist eine strukturierte, nachvollziehbare Begründung, dass ein System ausreichend sicher ist. 
Es verknüpft Sicherheits-Behauptungen (Claims) mit der Argumentationsstrategie und konkreter Evidenz 
(Tests, Analysen, Reviews), inklusive der Annahmen und Geltungsbereich. Dadurch wird nachvollziehbar, 
warum die Risiken als beherrscht gelten und welches Restrisiko bleibt. Für den Nachweis der Sicherheit 
ist es zentral, weil es Prüfern und Behörden eine transparente, prüfbare Kette von Anforderungen → Maßnahmen → Nachweisen liefert.

Die folgenden Punkte sind die in ISO/PAS~8800 geforderten Anforderungen um ein Assurance-Argument zu erstellen:

\begin{itemize}
  \item \textbf{(a) Definition des KI-Systems} (aus externen Quellen), einschließlich:
    \begin{enumerate}[label=\arabic*)]
      \item Spezifikation der dem KI-System zugewiesenen Sicherheitsanforderungen;
      \item Definition des technischen Kontexts innerhalb des übergeordneten Systems (z.\,B.\ Schnittstellen, Bedingungen/Trigger, unter denen die KI-Funktion aktiviert wird, usw.);
      \item Spezifikation des Eingaberaums.
    \end{enumerate}
  \item \textbf{(b) Anforderungen an das Assurance-Argument und die Arbeitsprodukte} für das KI-System (aus externen Quellen). Diese Anforderungen können aus dem Assurance-Argument des übergeordneten Systems sowie aus den Sicherheitsmanagement-Prozessen abgeleitet werden.
  \item \textbf{(c) Die Arbeitsprodukte des KI-Lebenszyklus}
\end{itemize}

\subsection{(9) Ableitung von KI-Sicherheitsanforderungen}
Es wird der Ablauf und Methoden beschrieben, wie Anforderungen abgeleitet werden können. 
Da diese sich nicht auf Systeme, in der eine KI eingebettet wird beziehen oder wie diese 
konkret auf KI Modelle angewendet werden können, ist diser Abschnitt sehr high-level. 
Ein essentielles Beispiel, wie Anforderungen verändert werden müssen ist folgendes:

\[
\begin{array}{c}
\text{Anforderung an}\\
\text{Wahrscheinlichkeitn}\\
\text{von Fehlern, die}\\
\text{selten auftreten}
\end{array}
\;\underset{\text{verfeinert zu}}{\longrightarrow}\;
\underbrace{\begin{array}{c}
\text{Anforderung an}\\
\text{relative Frequenzn}\\
\text{von Fehlern, die}\\
\text{selten auftreten}
\end{array}}_{\substack{\text{Quantitative Anforderung}}}
\;\land\;
\underbrace{\begin{array}{c}
\text{Anforderung an}\\
\text{zusätzliche technische}\\
\text{Maßnahmen zum}\\
\text{kontrollieren oder}\\
\text{reduzieren der}\\
\text{Unsicherheit bei}\\
\text{der Bewertung}\\
\text{der quantitativen}\\
\text{Anforderungen}\\
\end{array}}_{\text{Qualitative Anforderung}}
\]


\subsection{(10) Auswahl von KI-Technologien sowie Architektur- und Entwicklungsmaßnahmen}
Die Ziele dises Abschnittes sind bereits bei Inhalten beschrieben. Es wird beschrieben, dass Sicherheitsanforderungen 
durch zwei Maßnahmen erfüllt werden können: Die Architektur und die Entwicklung. Weiter können beim Designprozess 
nicht immer alle Anforderungen erfüllt werden, was dazuführt, dass diese angepasst werden müssen. Wichtig sind 
die generellen Anforderungen, nach denen die KI-Technologien sowie Architektur- und Entwicklungsmaßnahmen entschieden werden:

\begin{itemize}
  \item Rechtfertigung, dass die Sicherheitsanforderungen erfüllt sind.
  \item Sicherheitsanforderungen müssen den Komponenten zugewiesen sein.
  \item Architektur- und Entwicklungsmaßnahmen helfen vom Design her Fehler von KI-Komponenten zu verhindern um Sicherheitsanforderungen zu erfüllen und das Risiko zu minimieren.
  \item Die Effektivität gewählter Maßnahmen muss mittels Begründungen unterstützt werden.
  \item Die Sicherheitsanalyse des KI-System-Output zeigt, ob Anforderungen erfüllt werden können.
  \item Unterschiede zwischen Entwicklungsumgebung und Zielumgebung muss identifiziert werden.
  \item KI-Modelle müssen mittels validierungs Datensätzen trainiert und evaluiert werden.
\end{itemize}

Der letzte Punkt schließt indirekt RL Methoden aus, auch wenn diese, durch diese Norm allgemein nicht ausgeschlossen werden. 
Das stellt einen kleinen Widerspruch dar, was aufzeigt, dass auch Normen iterativ verbessert werden müssen. Anzumerken ist, dass 
dies eine erst Version ist und daher noch keine iterativen Schritte seit der Erstveröffentlichung durchlaufen hat.

\paragraph{Beispiel von Architektur- und Entwicklungsmaßnahmen}
Weiter befinden sich in diesem Abschnitt Beispiele wie Architektur- und Entwicklungsmaßnahmen aussehen können. 
Konkrete Maßnahmen bezüglich Robustheit, Zuverlässigkeit, etc. stützen auf bestpractise Methoden und 
liefern daher keine zusätzlichen Erkenntnisse wann eine KI sicher ist.

\subsection{(11) Datenbezogene Aspekte}
Die hier beschriebenen Aspekte beziehen sich hauptsächlich auf: un-/ semi- und supervised learning. 
Da diese sich ebenfalls auf bestpractise Methoden stützen, werden sie in dieser Arbeit nicht zusätzlich 
aufgeführt, da diese nur wenig zur Erkenntnissgewinnung beitragen. Diese Methoden sind allgemein beim 
Arbeiten mit großen Datenmengen und auch bei ML von essentieller Bedeutung. Sie sind weitverbreitete 
Standards verschiedener Industrien und haben auch Einfluss auf verschiedene Punkte dieser Arbeit.

\subsection{(12) Verifikation und Validierung des KI-Systems}
Diese Kapitel geht sehr ``highlevel'' auf Verifikation und Validierung ein. Unter anderem mit einigen Anforderungen 
an benötigten Tests und des Workflows mit wichtigen Stichpunkte wie, dass SOTIF Trigger Bedingungen beschreibt. 
Desweiteren wird auch eine Liste möglicher Testmethoden vorgeschlagen um Komponenten zu testen. Für diese 
Arbeit sind folgende Punkte wichtig:

\paragraph{KI spezifische Herausforderungen mit Erklärungen}
\begin{itemize}
  \item fehlen von präzisen Sicherheitsanforderungen
  \item vielfältige Eingaben aus verschiedenen Quellen (z.B. Kombination aus Radar und Kamera)
  \item DNNs haben komplexe Architekturen und werden heuristische Trainiert (basierend auf Erfahrungswerten)
  \item das Basieren auf großen Datensätzen
  \item nicht vorhersagbares fehlerhaftes Verhalten
  \item Einschränkung struktureller Abdeckung
  \item kleine Änderungen (Umwelt) bringen große Auswirkungen mit sich (input VS output)
  \item nicht ausreichende Datensätze führen zu ungewollten Verhalten
\end{itemize}

Weiter wird ein Beispiel zur Objekterkennung angeführt, dass sowohl ein Lidar als auch eine 
Kamera ein Objekt nicht erkennen können durch einen Absorptionseffekt eines Objektes. 
Um zu zeigen, dass diese nicht stochastisch unabhängig sind.

\subsection{(13) Sicherheitsanalyse von KI-Systemen}
Es werden recht allgemein Sicherheitsanalysen und Workflows beschrieben mit berücksichtigung von KI. 
Speziell angesprochene Punkte, sind: 
\begin{itemize}
  \item einige Datenspezifische Verweise
  \item die Unterscheidung von KI-Komponenten in KI-Modelle und nicht KI-Modelle (Steuerlogik)
  \item die am häufigsten auftretenden Fehlern bei KIs sind:
    \begin{itemize}
      \item Probleme bei Datenspezifikationen und dem sammeln von Daten
      \item Probleme im Design und Implementierung
      \item Probleme beim spezifizieren von Anforderungen
    \end{itemize}
  \item Herausforderungen und schwierige Eigenschaften bei Sicherheitsanalyse Techniken sind:
    \begin{itemize}
      \item Trainingsdaten statt Systemspezifikationen
      \item unklare Systemarchitekturen
      \item Unsicherheiten
      \item Erklärbarkeit
      \item nichtlineares Verhalten
      \item sich verändernde Umgebung
      \item komplexe Interaktionen und schwer verständliche Zusammenhänge
    \end{itemize}
\end{itemize}


\subsection{(15) Vertrauen in die Nutzung von KI-Entwicklungsframeworks und -Softwaretools}
In diesem Kapitel beschrieben sind bestpractise Methoden gestützt mit Beispielen und ein Fokus auf 
PFMEA (Potential Failure Modes and Affects Analysis). 


\section{Annex}
Im Anhang befinden sich wichtige Beispiele und Zusatzinformationen. Folgend sind die 
für diese Arbeit wichtigsten Punkte:
\paragraph{Annex - C; ISO 26262 Lückenanalyse für ML}
Für den Einsatz von ML wird Hardware von einem oder mehreren CPUs empfohlen und 
für Neuronale Netze GPUs. \\
Vorallem bei der Verifikation, Analyse, etc. wird auf Simulationen und Heuristische 
bzw. Experten Aswertung verwiesen.\\
Die Anforderung an die Richtigkeit der Ausführung des Codes und Datenfluss muss entfernt 
werden, was auf den Probabilistischen Output einer KI zurückzuführen ist. Das bedeutet, 
dass der Code nicht mehr deterministisch sein kann und eine MC/DC Coverage (Modified Condition/Decision Coverage) 
nicht mehr möglich ist. (MC/DC ist eine Anforderung an die Eindeutigkeit des Codes; Jeder Ja/Nein Input oder 
Bedingung hat einen Eindeutigen Output.)
\paragraph{Annex - D; Detaillierte Betrachtungen zu sicherheitsrelevanten Eigenschaften von KI-Systemen}
Hier werden Sicherheits bezogene Eigenschafte beschrieben und zugeordnet, diese hier wiederzugeben 
wäre exessiv und trägt nicht unmittelbar dem Verständnis bei.
\paragraph{Annex - G; Architektur- und Entwicklungsmaßnahmen für KI-Systeme}
Hier werden Detailbeschreibungen angeführt zu folgenden Themen, welche jeweils untergliedert sind:
\begin{itemize}
  \item Beispiele für Architektur und Entwicklungsmaßnahmen für KI
  \item Qualitative und qantitative Analyse von KI Architekturen
  \item Datenverteilungen und ihr Einfluss auf KI Modelle
  \item Trainingssicherheitsmaßnahmen
  \item Überwachung und KI System Modifikationen
  \item Ausrichtung der Absichten
  \item Berücksichtigungen im Zusammenhang mit der Zielausführungsumgebung
\end{itemize}


% -------------------------
% CHAPTER 6: Coverage-Analyse entlang ISO/PAS 8800
% -------------------------
\chapter{Coverage-Analyse}\label{chap:coverage}

\section{Vorgehen der Analyse}\label{sec:analyse-vorgehen}
Die Analyse fokussiert nicht die Wirksamkeit einzelner Techniken, sondern welche Ansätze zur Absicherung gefordert werden
(u.\,a.\ aus ISO/PAS~8800 sowie konzeptionellen Leitlinien) und in welchem Umfang Abdeckung erkennbar ist.
Als Anker dienen die in Kapitel~\ref{chap:anforderungen} definierten Anforderungsfelder. 




\newcommand{\covV}{\textcolor{green!40!black}{\textbf{V}}} % voll
\newcommand{\covT}{\textcolor{orange!80!black}{\textbf{T}}} % teilweise
\newcommand{\covK}{\textcolor{red!70!black}{\textbf{K}}}    % keine
\newcolumntype{L}[1]{>{\raggedright\arraybackslash}p{#1}}

\section{Vergleich: Luftfahrt-Forderungen vs.\ ISO 8800 (Coverage-Matrix)}
\setlength{\LTpost}{0pt}

\begin{longtable}{L{0.12\linewidth}L{0.4\linewidth}L{0.18\linewidth}L{0.12\linewidth}}
\caption{Gegenüberstellung der Forderungen (F-Avi-01…12) mit ISO/PAS 8800:2024 \textbf{Kapiteln}. Abdeckung: \covV{}=voll, \covT{}=teilweise, \covK{}=keine.}\\
\toprule
\textbf{ID} & \textbf{Kurzbezeichnung} & \textbf{ISO 8800 (Kapitel)} & \textbf{Abdeckung} \\
\midrule
\endfirsthead
\toprule
\textbf{ID} & \textbf{Kurzbezeichnung} & \textbf{ISO 8800 (Kapitel)} & \textbf{Abdeckung} \\
\midrule
\endhead
\midrule
\multicolumn{4}{r}{Fortsetzung auf nächster Seite}\\
\midrule
\endfoot
\bottomrule
\endlastfoot

F-Avi-01 & Human-centric, Level \& ConOps/HAT & 6, 7 & \covT \\[2pt]

F-Avi-02 & OD/ODD \& Abstain/Degradation & 9, 12, 14 & \covV \\[2pt]

F-Avi-03 & Knowledge \& Data Governance (ODD→Daten) & 11, 9, 8 & \covV \\[2pt]

F-Avi-04 & AI Assurance (Learning \& Explainability) & 8, 12, 15 & \covT \\[2pt]

F-Avi-05 & Betriebs-Recording & 14 & \covV \\[2pt]

F-Avi-06 & OOD \& Unvorhersagbarkeit& 9, 12, 14 & \covT \\[2pt]

F-Avi-07 & Generalisierung, Robustheit \& Stabilität & 9, 12 & \covV \\[2pt]

F-Avi-08 & Lebenszyklus \& W-Shape (Gates/Artefakte) & 7, 8, 14 & \covT \\[2pt]

F-Avi-09 & Kein Online-Lernen im Betrieb & 14, 10 & \covT \\[2pt]

F-Avi-10 & Einbettung in Safety-Systeme/Plattform & 10, 6 & \covV \\[2pt]

F-Avi-11 & Rare-Events-Management (Taxonomie/Katalog) & 11, 12, 9 & \covT \\[2pt]

F-Avi-12 & Use-Case-Verankerung (KPIs/Übergabe in Betrieb) & 6, 8, 12, 14 & \covT \\

\end{longtable}




\section{Diskussion der Abdeckung (EASA \& ISO 8800)}

\subsection{Überblick}
Die Gegenüberstellung zeigt eine solide Abdeckung zentraler Lebenszyklus- und Technikthemen durch die ISO/PAS~8800. 
Insbesondere Datenführung, V\&V sowie Betriebsmaßnahmen sind deutlich adressiert. 
Demgegenüber bleiben Aspekte mit starkem Human-Fokus (HAT/ConOps), genaue Ausprägungen von Erklärbarkeit, 
explizite Runtime-Indikatoren und organisatorische Governance-Elemente nur teilweise oder rahmenhaft geregelt.

\subsection{Stärken (voll abgedeckt, \covV)}
\begin{itemize}
  \item \textbf{F-Avi-02 OD/ODD \& Abstain/Degradation} (9, 12, 14): \\
  ISO verknüpft ODD/Unsicherheit mit Ableitung von Anforderungen, Test/Validierung und Maßnahmen im Betrieb. 
  Domänenabgrenzungen und der Umgang mit Degradation wird methodisch behandelt, gilt es aber projektspezifisch auszuarbeiten.
  \item \textbf{F-Avi-03 Knowledge \& Data Governance} (11, 9, 8): \\
  Vollständiger Dataset-Lebenszyklus mit Safety-Analyse, Validierung, Pflege und Rückführung in den Assurance-Nachweis. 
  \item \textbf{F-Avi-05 Betriebs-Recording} (14): Feld-Datenerhebung, Nutzung für Monitoring, Re-Bewertung/Änderungen sind fest verankert.
  \item \textbf{F-Avi-07 Generalisierung, Robustheit \& Stabilität} (9, 12): \\
  Prüfziele, Methoden und Bewertungslogik für Performance/Robustheit sind vorhanden (inkl.\ Testebenen und Integrationssicht).
  \item \textbf{F-Avi-10 Einbettung in Safety-Systeme/Plattform} (10, 6): \\
  Architektur- und Entwicklungsmaßnahmen, sowie Systemkontext und Einbettung werden adressiert.
\end{itemize}

\subsection{Teilabdeckungen (Lücken, \covT)}
\begin{itemize}
  \item \textbf{F-Avi-01 Human-centric, Level \& ConOps/HAT} (6, 7): \\
  Rollen, Interaktion und Management sind nur implizit vorhanden. Der Mensch als Zentrale Rolle 
  und konkrete Aufgabenverteilungen sind nicht explizit vorgegeben. Diese sind in den Anforderungen 
  oder Domänenabgrenzungen zu betrachten. Der Mensch wird bezüglich des Automatisierungsgrad berücksichtigt 
  und soll hier in die Anforderungen einfließen, wird aber nicht in den Mittelpunkt gestellt.
  \item \textbf{F-Avi-04 AI Assurance (Learning \& Explainability)} (8, 12, 15): \\
  Assurance-Argumente sind vollständig dargestellt. Erklärbarkeit wird ausdrücklich 
  gefordert, bleibt aber offen und ist projektspezifisch zu betrachten. Diese sind sehr offen 
  und nur beispielhaft konkretisiert.
  \item \textbf{F-Avi-06 OOD \& Unvorhersagbarkeit} (9, 12, 14): \\
  OOD/Unsicherheiten werden definiert und im Test/Betrieb behandelt. 
  Spezifische Runtime-Monitore/Trigger sind gefordert zu dokumentieren aber als Aufgabe der 
  SOTIF und mit Hinweisen zum ``living safety-case'' vermerkt. In der Luftfahrt werden z.B. 
  in den Anforderungen explizite Update Zeitpunkte festgehalten, im Automotive werden diese 
  unter anderem gesondert betrachtet.
  \item \textbf{F-Avi-08 Lebenszyklus \& W-Shape} (7, 8, 14): \\
  Ein vollständiger AI-Safety-Lebenszyklus ist vorhanden; das EASA-spezifische W-Shape (Gates/Artefakte) 
  als formale Struktur ist nicht Teil der ISO. Statdessen wird die klassische V-Shape als Struktur verwendet.
  \item \textbf{F-Avi-09 Kein Online-Lernen im Betrieb} (14, 10): \\
  ISO fordert Re-Training/Re-Validation/Re-Approval bei Änderungen, verbietet Online-Lernen jedoch nicht explizit. 
  Die EASA hingegen rät vom Online-Lernen stark ab.
  \item \textbf{F-Avi-11 Rare-Events-Management} (11, 12, 9): \\
  Edge-Case-Handling ist vorgesehen; ein verpflichtender Taxonomie-/Katalog-Mechanismus ist nicht vorgegeben. 
  Edge-cases sind fest zu halten, zu definieren, wie damit umzugehen ist, etc.. Sie sind aber projektspezifisch, 
  daher wird ein allgemeiner Katalog nicht gefordert und ist für jedes Projekt einzeln anzufertigen.
  \item \textbf{F-Avi-12 Use-Case-Verankerung} (6, 8, 12, 14): \\
  Kontext, KPIs, Validierung und Betriebsübergang sind vorhanden; projektinterne Use-Case-Artefakte/KPIs 
  müssen konkretisiert werden. Diese sind hier zwar gefordert, es wird aber darauf hingewiesen, 
  dass diese Projektabhängig sind.
\end{itemize}

\subsection{Interpretationsrahmen}
Die ISO/PAS~8800 setzt auf prinzipienbasierte Anforderungen: Sie definiert Ziele, Workflows und Work Products, lässt aber 
Gestaltungsspielraum für projektspezifische Ausprägungen (z.\,B.\ HAT/HMI, Monitore, interne Gates). 
Die ausgewiesenen Teilabdeckungen (\covT) sind daher Erwartungsräume, die im Projektkontext präzisiert werden sollten, 
ohne der ISO zu widersprechen.

\subsection{Zwischenfazit}
Der Abgleich der EASA-Dokumente mit der ISO/PAS~8800 zeigt, dass die in dieser Arbeit behandelten Aspekte von der Norm
entweder vollständig oder zumindest teilweise abgedeckt werden. Die ISO/PAS~8800 ist bewusst
prinzipienbasiert angelegt und lässt projektspezifischen Interpretations- und Gestaltungsspielraum. Zugleich treten
Branchenunterschiede zutage: Die Luftfahrt verankert unter anderem Anforderungen und Einsatzgrenzen (OD/ODD) stark im
Prozess und verfolgt eine deutlich mensch-zentrierte Auslegung, während die Automobilbranche viele Fragestellungen auch
aus der SOTIF-Perspektive adressiert und den Menschen zwar berücksichtigt, ihn jedoch nicht in gleicher Weise ins Zentrum
stellt.

Eine konkrete Übertragbarkeit der Anpassungen der ISO 26262 im vergleich zur DO 178C gilt es noch zu untersuchen. 
Dies würde den Umfang dieser Arbeit überschreiten, ist aber notwendig um die Anwendbarkeit der hier identifizierten 
Abdeckung beurteilen zu können.